% ============================================================
% CHAPTER 8: REQUIREMENTS ELICITATION
% MCDA 5585 / 5586 Major Project Report
% Bhavik Kantilal Bhagat | A00494758 | Saint Mary's University
% ============================================================

\chapter{Requirements Elicitation}
\label{chap:requirements}

During my internship I worked with a small set of internal stakeholders
whose needs shaped the functional and non-functional requirements of the
IA-IT-SRE Infrastructure monitoring platform.

% ────────────────────────────────────────────────────────────
\section{Stakeholder Identification}
\label{sec:stakeholders}
% ────────────────────────────────────────────────────────────

Table~\ref{tab:stakeholders} enumerates the key
stakeholders, their roles, and their specific interests in the
project.

\begin{table}[ht]
\centering
\caption{Stakeholder Identification}
\label{tab:stakeholders}
\rowcolors{2}{tablerowA}{tablerowB}
\begin{tabularx}{\textwidth}{p{2.5cm}p{3.8cm}X}
\hline
\rowcolor{smubrand}\textcolor{white}{\textbf{Stakeholder}} & \textcolor{white}{\textbf{Role}} & \textcolor{white}{\textbf{Interest / Influence}} \\
\hline
Kishor & Manager, CTM; Industry Supervisor & Primary sponsor; set engineering priorities; escalated IISS-487 to Critical; final acceptance authority \\
Horace & IT Senior Developer, RPA Platform Lead & RPA subject matter expert; delivered Blue Prism and UiPath KT; defined RPA monitoring scope \\
Kri & IT Developer, SRE Engineer & Delivered golden signals and alarm threshold training; reviewed dashboard design \\
IA IT SRE Team & Internal SRE consumers & Primary users of the Grafana dashboard and CloudWatch alarm notifications \\
Meridian Team & Kafka/ECS architecture owners & Provided MSK consumer group structure, ECS cluster topology, and OpenSearch integration context \\
Sridhar \& Soundarya & Co-interns & Parallel development of related monitoring features; design review and test validation \\
\hline
\end{tabularx}
\end{table}

% ────────────────────────────────────────────────────────────
\section{Requirements Elicitation}
\label{sec:requirements_elicitation}
% ────────────────────────────────────────────────────────────

The elicitation methods used were appropriate for an
in-house engineering context: structured knowledge-transfer sessions,
incident post-mortems, and JIRA story refinement, rather than formal
user interviews or focus groups.

\subsection{Elicitation Methods}

\paragraph{Knowledge-Transfer Sessions (Training Week, Apr 13--18, 2026).}
The first week of full-time engagement was structured as an intensive
knowledge-transfer programme covering the full breadth of the IA
platform. Sessions with Horace on the Blue Prism and UiPath
architectures, with Bhanu Teja Murari on CitcoWorks and the Event Store
platform, and with the Meridian team on the Kafka event pipeline
(comprising the Event Routing Service, Action Processor, and OpenSearch
consumer) provided the architectural grounding needed to define
meaningful SLIs. Key requirements insight from training week: the
\AE xeo Treasury platform (cloud-based wire approval and fund movement)
and the MESO platform (high-memory autoscale issue encountered later) each
had distinct resource profiles requiring separate monitoring strategies.

\paragraph{Incident Post-Mortem (Apr 17, 2026).}
The Meridian production incident yielded the single most important
requirements insight of the entire engagement: MSK consumer
lag is the primary leading indicator of Meridian health. The
post-mortem made concrete the latency between signal onset and
detection under the reactive model, and directly specified the target
for the lag alarm evaluation period (NFR2: $\leq$1~minute).

\paragraph{Golden Signals Session (May 25, 2026).}
Kri's dedicated SRE methodology session translated the theoretical
four golden signals into concrete metric choices for the IA platform.
This session produced the initial SLI/SLO target table and established
the design constraint that all metrics must be queryable through
Amazon CloudWatch --- ruling out external time-series databases such as
Prometheus and anchoring the architecture to CloudWatch Metric Insights
as the sole data source.

\paragraph{JIRA Story Refinement.}
The dashboard project went through multiple refinement cycles between May~6 and
May~11 (infrastructure deployment date). These sessions with Kishor
and Kri surfaced the infrastructure-as-code requirement (FR7), the
multi-environment reuse requirement (NFR5), and the credential
management constraint (NFR7). The elevation of the dashboard to Critical
priority on May~6 also functionally set NFR1 (dashboard refresh
latency $\leq$30~seconds) as a hard constraint rather than an aspirational
target.

\subsection{Key Insights}

Three insights from the elicitation process materially shaped the
resulting design:

\begin{enumerate}
    \item \textbf{Tag-based resource discovery is non-trivial at Citco
          scale.} The training week revealed that the \texttt{citcoworks}
          and \texttt{meridian} environments share the
          \textbf{Shared Cloud Artifacts} BudgetCode tag, meaning
          a simple BudgetCode lookup conflates resources from
          two operationally distinct environments. This directly
          motivated FR9 (tag-based discovery with AppName
          fallback).

    \item \textbf{The primary audience for the dashboard is the SRE
          team, not business operators.} Business operations teams would
          continue to use their existing domain-specific tools; the
          dashboard needed to be useful for the SRE first-responder,
          not designed as a business reporting tool. This shaped the
          panel selection (technical metrics over business KPIs) and
          the deep-link design (panels link directly to the relevant AWS
          Console page).

    \item \textbf{Infrastructure-as-Code is a first-class requirement,
          not an after-thought.} The team had experienced the cost of
          manually provisioned monitoring infrastructure in previous
          projects (configuration drift, undocumented alarms). The JIRA
          refinement sessions made FR7 an explicit acceptance criterion
          rather than a \textbf{nice-to-have}.
\end{enumerate}

% ────────────────────────────────────────────────────────────
\section{Functional and Non-Functional Requirements}
\label{sec:requirements}
% ────────────────────────────────────────────────────────────

The following requirements were derived from the elicitation activities
described above. They governed the design and acceptance criteria for
all six work streams of my internship.

\subsection{Functional Requirements}

Table~\ref{tab:functional_requirements} enumerates the nine functional
requirements, each with the primary deliverable that satisfies it.

\begin{table}[ht]
\centering
\caption{Functional Requirements}
\label{tab:functional_requirements}
\rowcolors{2}{tablerowA}{tablerowB}
\begin{tabularx}{\textwidth}{p{0.9cm}Xp{5.0cm}}
\hline
\rowcolor{smubrand}\textcolor{white}{\textbf{ID}} & \textcolor{white}{\textbf{Requirement}} & \textcolor{white}{\textbf{Satisfied By}} \\
\hline
FR1 & Dashboard displays real-time Lambda metrics (invocations, errors, duration, throttles) per platform & Grafana AMG, Lambda rows (IISS-487) \\
FR2 & Dashboard displays ECS CPU and memory utilization per cluster and service & Container Insights + ECS rows (IISS-487, IISS-483) \\
FR3 & Dashboard displays MSK consumer lag per consumer group & MSK rows; \texttt{EstimatedMaxTimeLag} alarm (IISS-487) \\
FR4 & Dashboard displays SQS queue depth and oldest message age & SQS rows in Grafana; SQS monitoring (IISS-482) \\
FR5 & CloudWatch alarms provisioned for all monitored resource types with SNS routing to SRE team & Dynamic alarms via CloudFormation (IISS-487) \\
FR6 & Dashboard supports variable-based filtering by platform and resource & Grafana template variables (IISS-487) \\
FR7 & All monitoring infrastructure deployed as Infrastructure-as-Code & 4 CloudFormation stacks; Lambda-backed IaC provisioner (IISS-487) \\
FR8 & Inventory API enumerates Lambda, ECS, EC2, and S3 resources per platform via REST & Flask REST API, discoverer registry (IISS-507) \\
FR9 & Inventory API supports two-stage tag-based discovery for shared BudgetCode platforms & Two-stage AppName filter in \texttt{citcoworks}/\texttt{meridian} discoverers (IISS-507) \\
\hline
\end{tabularx}
\end{table}

\subsection{Non-Functional Requirements}

Table~\ref{tab:nonfunctional_requirements} enumerates the seven
non-functional requirements, each with the measurable target and the
achieved value where applicable.

\begin{table}[ht]
\centering
\caption{Non-Functional Requirements}
\label{tab:nonfunctional_requirements}
\rowcolors{2}{tablerowA}{tablerowB}
\begin{tabularx}{\textwidth}{p{0.9cm}Xp{2.2cm}p{3.0cm}}
\hline
\rowcolor{smubrand}\textcolor{white}{\textbf{ID}} & \textcolor{white}{\textbf{Requirement}} & \textcolor{white}{\textbf{Target}} & \textcolor{white}{\textbf{Achieved}} \\
\hline
NFR1 & Dashboard refresh enables sub-minute incident detection & $\leq$30~sec & $\leq$30~sec \\
NFR2 & CloudWatch alarm evaluation period for MSK lag and Lambda error rate & $\leq$1~min & 1~min \\
NFR3 & Inventory API response time for full resource enumeration & $\leq$10~sec & $<$10~sec \\
NFR4 & Inventory API automated test coverage & $\geq$95\% & 99\% \\
NFR5 & CloudFormation templates parameterized for multi-environment reuse & Parameterized YAML & Achieved \\
NFR6 & Dashboard functional after Grafana version upgrades without manual rework & Standard plugin queries & Achieved (10.4$\to$12.4) \\
NFR7 & Zero hardcoded credentials --- all secrets via AWS Secrets Manager & No plaintext secrets & Achieved \\
\hline
\end{tabularx}
\end{table}

\subsection{Requirements Traceability}

All nine functional requirements and all seven non-functional
requirements were satisfied by my delivered work streams. The
test suite for the inventory API provides quantitative evidence for NFR3
and NFR4. NFR6 was validated in practice when the Grafana workspace
was upgraded from version 10.4 to version 12.4 on May~13, 2026 ---
a major version migration that required updating CloudWatch Metric
Insights query syntax across all panels but did not require any
architectural change to the underlying monitoring design.

The most technically demanding requirement was FR9, which exposed a
fundamental ambiguity in the AWS tagging strategy for the IA platform.
My discovery that \texttt{citcoworks} and \texttt{meridian} share the
same BudgetCode tag value (\textbf{Shared Cloud Artifacts})
was not documented anywhere before I built the inventory API. The
solution --- designating AppName as a mandatory secondary
discriminator for these two environments --- became the standard pattern
for all resource discovery going forward and was documented in the
inventory API's README as a platform-wide architectural note.

NFR7 received an unexpected real-world validation during the
CAIS Deadline hotfix on Jun~8, 2026. The incident revealed
that the ECS container's credential loading pattern (reading from
Secrets Manager once at module import time via \texttt{constants.py})
was fragile under secret rotation. The hotfix I delivered introduced fail-fast
credential validation --- a design improvement that reinforced NFR7
by ensuring that absent or expired credentials from Secrets Manager
produce an immediate, visible failure rather than a silent downstream
authentication error.

\subsection{Out-of-Scope Requirements}

The following were evaluated and explicitly descoped:

\begin{itemize}
    \item \textbf{Automated error budget burn-rate alerting.} Computing
          burn rate requires Metric Math expressions that divide the
          rolling error count by the rolling SLO budget. While the SLO
          targets and SLI metrics are now in place, configuring
          burn-rate alarms requires additional CloudFormation work that
          I descoped from the dashboard project.

    \item \textbf{Prometheus as a data source.} A self-managed Prometheus
          server with a CloudWatch Exporter sidecar was evaluated as an
          alternative data layer. It was ruled out because Amazon Managed
          Grafana does not support a self-managed Prometheus endpoint
          without Amazon Managed Service for Prometheus (AMP), which
          would add infrastructure cost and operational overhead.
          CloudWatch Metric Insights SQL provides equivalent
          multi-resource aggregation capability without the additional
          stack. Documented as Decision~3 in
          Section~\ref{sec:decision_cloudwatch_only}.

    \item \textbf{Pyroscope (continuous profiling).} Pyroscope is a
          Grafana-ecosystem continuous profiling tool that would surface
          CPU flame graphs and memory allocation traces at the function
          level. While it would complement the golden signals dashboard,
          it requires an instrumentation agent deployed alongside each
          monitored service. Given the scope of the current engagement
          and the constraint that the Grafana workspace uses only
          CloudWatch as its data source, Pyroscope was identified as a
          post-internship enhancement.
\end{itemize}


% Keep all floats within this chapter
\FloatBarrier
